% !TeX root = ../0_system_design_report.tex
% Figure 1 -- upgraded system architecture.  Portrait, sized to fill a page.
% The point the drawing must make: machine protection (Interface Chassis) and
% personnel safety (PPS Interface Box) are architecturally separate paths.
%
% Layout rule that keeps every line clear: the lower row is offset by half a
% pitch from the upper row, so each lower-row box sits under a GAP in the upper
% row and each upper-row box sits over a GAP in the lower row.  Every drop is
% therefore straight -- the Channel Access feeds reach the lower row without
% crossing an upper-row box, and the upper row reaches the Interface Chassis
% without crossing a lower-row box.  Only the HVPS PLC is routed sideways,
% down the clear lane to the right of the heater controller.
\begin{tikzpicture}[font=\scriptsize, x=1mm, y=1mm]

% ===== operator and expert layers =========================================
\node[ctrl,text width=54mm,anchor=north] (op) at (32,0)
  {\textbf{OPERATOR LAYER}\\
   EDM panels $\cdot$ web dashboard\\ EPICS archiver $\cdot$ logging};
\node[ctrl,text width=76mm,anchor=north] (ex) at (114,0)
  {\textbf{EXPERT LAYER}\\
   MATLAB tools (Dimtel LLRF9 + in-house)\\
   expert EPICS panels $\cdot$ system diagnostics};

% ===== EPICS coordinator ==================================================
\node[ctrl,text width=154mm,anchor=north] (epics) at (80,-24)
  {\textbf{EPICS (soft IOC) COORDINATOR}\\
   state machine $\cdot$ HVPS loop $\cdot$ tuner manager $\cdot$ fault manager
   $\cdot$ diagnostics $\cdot$ heartbeat monitoring};
\draw[sig] (op.south) -- (op.south|-epics.north);
\draw[sig] (ex.south) -- (ex.south|-epics.north);

\def\yCA{-44}
\coordinate (caY) at (0,\yCA);
\draw[bus] (10,\yCA) -- (126,\yCA);
% the rail label sits ABOVE the rail, clear of the coordinator drop at x = 80
\node[tag,anchor=south west,fill=none] at (12,-42.6)
  {EPICS Channel Access\quad($\sim$\qty{1}{\hertz} supervisory)};
\draw[sig] (epics.south) -- (80,\yCA);

% ===== upper subsystem row ================================================
% centres 16 / 50 / 84 / 118, width 26  ->  spans 3-29, 37-63, 71-97, 105-131
\def\yU{-58}
\node[rf,text width=26mm,anchor=west]  (l1)  at (3,\yU)
  {\textbf{LLRF9 \#1}\\ field control\\ tuner control\\ interlocks};
\node[rf,text width=26mm,anchor=west]  (l2)  at (37,\yU)
  {\textbf{LLRF9 \#2}\\ field control\\ tuner control\\ interlocks};
\node[blk,text width=26mm,anchor=west] (mot) at (71,\yU)
  {\textbf{MOTOR CTRL}\\ Galil DMC-4143\\ Ethernet,\\ heartbeat};
\node[plc,text width=26mm,anchor=west] (hp)  at (105,\yU)
  {\textbf{HVPS PLC}\\ CompactLogix\\ voltage control};
\foreach \n in {l1,l2,mot,hp}{\draw[sig] (\n.north|-caY) -- (\n.north);}
% The two LLRF9 interlock chains are wired directly to each other, not via the
% Interface Chassis.  Unlabelled on purpose: the only clear space for a label
% is the lane the Waveform Buffer feed uses.  Explained in the caption.
\draw[{Stealth[length=2mm]}-{Stealth[length=2mm]},thick] (l1.east) -- (l2.west);

% ===== lower subsystem row ================================================
% centres 33 / 67 / 101 / 135, width 26  ->  spans 20-46, 54-80, 88-114, 122-148
\def\yL{-84}
\node[blk,text width=26mm,anchor=west] (wb)   at (20,\yL)
  {\textbf{WAVEFORM}\\\textbf{BUFFER SYSTEM}};
\node[blk,text width=26mm,anchor=west] (arc)  at (54,\yL)
  {\textbf{ARC DETECT}\\ AFT ARC4};
\node[plc,text width=26mm,anchor=west] (mp)   at (88,\yL)
  {\textbf{RF MPS PLC}\\ ControlLogix 1756\\ collector-power calc};
\node[blk,text width=26mm,anchor=west] (heat) at (122,\yL)
  {\textbf{HEATER CTRL}\\ programmable AC supply};
% straight drops: 33, 67 and 101 all fall in gaps of the upper row
\foreach \n in {wb,arc,mp}{\draw[sig] (\n.north|-caY) -- (\n.north);}
\draw[sig] (mp.east) -- (heat.west);

% ===== interface chassis ==================================================
\node[plc,text width=154mm,anchor=north] (ic) at (80,-104)
  {\textbf{INTERFACE CHASSIS}\quad(new)\\
   first-fault detection $\cdot$ optocoupler / galvanic isolation $\cdot$
   fiber I/O $\cdot$ AND-gate permit logic};

% straight drops: 16, 50, 84 all fall in gaps of the lower row
\foreach \n in {l1,l2,mot}{\draw[sig] (\n.south) -- (\n.south|-ic.north);}
\foreach \n in {wb,arc,mp}{\draw[sig] (\n.south) -- (\n.south|-ic.north);}
% HVPS PLC takes the clear lane to the RIGHT of the heater controller
\draw[draw=cNeut,thick] (hp.south) |- (152,-72);
\draw[sig] (152,-72) -- (152,-104);

% ===== external permits ===================================================
\node[blk,text width=60mm,anchor=north] (ext) at (38,-130)
  {\textbf{MACHINE PROTECTION SAFETY SYSTEMS}\\
   SPEAR MPS $\cdot$ orbit interlock $\cdot$ external permits\\
   {\tiny feed into the Interface Chassis}};
\draw[sig] (ext.north) -- (ext.north|-ic.south);

% ===== HVPS power section =================================================
\node[pwr,text width=60mm,anchor=north] (hvps) at (116,-132)
  {\textbf{HVPS (power section)}\\
   transformer, rectifier, crowbar, thyristor stacks,\\
   grounding tank, Ross switch, filter inductors};
% leaves the BOTTOM edge of the Interface Chassis, not its top
\draw[fib] (ic.south -| hvps.north) -- (hvps.north);
\node[tag,anchor=west,align=left] at (119,-123)
  {3$\times$ fiber optic\\ SCR ENABLE $\cdot$ CROWBAR $\cdot$ STATUS\\
   {\tiny bypasses the PLC}};

% ===== PPS chain (separate path) ==========================================
\node[blk,text width=60mm,anchor=north,draw=cPwr!70,fill=cPwr!5] (ppsbox) at (116,-158)
  {\textbf{PPS INTERFACE BOX}\\
   Bud enclosure $\cdot$ 4 relays $\cdot$ status LEDs\\
   lockable connector\\
   HVPS contactor and Ross switch control};
\draw[sig] (hvps.south) -- (hvps.south|-ppsbox.north);
\node[tag,anchor=west,align=left] at (119,-151)
  {PPS interlock signals\\ {\tiny HVPS contactor, Ross switch}};

\node[blk,text width=60mm,anchor=north,draw=cPwr!70,fill=cPwr!5] (pps) at (116,-184)
  {\textbf{SPEAR PPS}\\ Personnel Protection System};
\draw[sig] (ppsbox.south) -- node[right,tag]{PPS chain signals} (pps.north);
\end{tikzpicture}
